\begin{proof}[Proof of \cref{obj:lem:selected-mesh}]
Throughout the proof, write
\[
m=\lceil\beta\rceil-1,\qquad
\mathcal I_m=\{\alpha\in\mathbb N_0^d:|\alpha|\le m\},
\qquad
p_{\mathrm{coef}}=|\mathcal I_m|,
\qquad
J=(m+1)^d,
\]
and use the template width \(\eta\) and least eigenvalue \(\lambda_0\) from
\cref{obj:def:template}. Also write
\[
D_\gamma=\frac{d\gamma}{\gamma-1},
\qquad
r_{\gamma,n}=h_{\gamma,n}^{\beta},
\qquad
\log_+(u)=\max\{0,\log u\}\quad(u>0).
\]

\begin{enumerate}
\item Membership in \(\mathcal P_\gamma\), as specified by
\cref{obj:def:model}, supplies the treated response bound in
\cref{obj:ass:bounded-potential-outcomes}. In particular, with
\[
\overline y_P(x)=\mathbb E_P[Y\mid A=1,X=x],
\]
the treated conditional law satisfies, for \(P_X\)-almost every \(x\),
\[
|\overline y_P(x)|\le B,\qquad
\mathbb E_P\!\left[
 e^{t(Y-\overline y_P(x))}\mid A=1,X=x
\right]\le e^{B^2t^2/2}
\quad(t\in\mathbb R),
\]
and every displayed exponential moment is integrable. These are the
outcome properties used below.
% lean: CausalSmith.Stat.WeakOverlap.selectedMesh_envelope

\item We first construct the uniform count bound. Define the positive mass
coefficient
\[
c_{\mathrm{mass}}(\gamma)
=
\frac{\gamma-1}{\gamma}\,
C^{-1/(\gamma-1)}
(c_f\eta^d)^{\gamma/(\gamma-1)},
\qquad
\kappa_*=\min_{\gamma\in\Gamma}c_{\mathrm{mass}}(\gamma)>0.
\]
The minimum exists and is positive because the displayed function is
continuous and positive on the compact interval \(\Gamma\).
Every template microcell at mesh \(h\) has volume \(\eta^dh^d\).
Consequently, \cref{obj:ass:density} and the setwise bound in
\cref{obj:lem:ordered-mass} give
\[
q_{Q\ell}\ge c_{\mathrm{mass}}(\gamma)h^{D_\gamma}
\ge\kappa_*h^{D_\gamma}.
\]

Choose a fixed enlargement \(A_{\mathrm{mesh}}\ge1\) such that
\[
4\le
\kappa_*A_{\mathrm{mesh}}^{\,2\beta+D_{\gamma_{\max}}}.
\]
Such a choice exists because \(2\beta+D_{\gamma_{\max}}>0\).
Since
\[
D_\gamma=d+\frac d{\gamma-1},
\]
the effective dimension decreases with \(\gamma\), and hence
\[
4\le
\kappa_*A_{\mathrm{mesh}}^{\,2\beta+D_\gamma}
\qquad(\gamma\in\Gamma).
\]

The candidate collection in \cref{obj:synth_3} has largest dyadic index
\[
j_{\max,n}
=
\left\lfloor\frac{\log_2 n}{2\beta+d}\right\rfloor.
\]
The floor inequality yields
\[
2^{-j_{\max,n}}\le 2n^{-1/(2\beta+d)}.
\]
Moreover,
\[
h_{\gamma_{\max},n}\le h_{\gamma,n}
\le h_{\gamma_{\min},n},
\qquad
\frac1{2\beta+d}-\frac1{2\beta+D_{\gamma_{\max}}}>0.
\]
The positive power in the last display tends to infinity when applied to
\(n\). Thus, for all sufficiently large \(n\),
\[
2^{-j_{\max,n}}
\le A_{\mathrm{mesh}}h_{\gamma_{\max},n}
\le A_{\mathrm{mesh}}h_{\gamma,n},
\qquad
A_{\mathrm{mesh}}h_{\gamma,n}
\le A_{\mathrm{mesh}}h_{\gamma_{\min},n}\le1.
\]
Rounding upward within the finite dyadic collection therefore gives a
candidate \(h_{\mathrm{o}}\) satisfying
\[
A_{\mathrm{mesh}}h_{\gamma,n}
\le h_{\mathrm{o}}
\le2A_{\mathrm{mesh}}h_{\gamma,n}.
\]
Indeed, successive candidate widths differ by a factor of two; the
smallest candidate width at least
\(A_{\mathrm{mesh}}h_{\gamma,n}\) satisfies these inequalities, including
when either endpoint is itself a candidate.

For the simultaneous concentration budget, define
\[
M_{\mathrm{count},n}
=
J\sum_{j=0}^{j_{\max,n}}2^{jd},
\qquad
s_n=\log\!\left(2(M_{\mathrm{count},n}+1)n^2\right).
\]
Here \(2^{jd}\) counts the partition cubes and \(J\) counts the microcells
within each cube. Since \(2\beta+d\ge1\),
\[
2^{j_{\max,n}}\le n,\qquad
j_{\max,n}\le2^{j_{\max,n}}\le n,
\]
and therefore
\[
M_{\mathrm{count},n}
\le(j_{\max,n}+1)(2^{j_{\max,n}})^dJ
\le(n+1)n^dJ.
\]
For \(n\ge2\), this implies
\[
2(M_{\mathrm{count},n}+1)n^2
\le8(J+1)n^{d+3},
\qquad
s_n\le(d+3)\log n+\log(8(J+1)).
\]

To make the count threshold dominate this budget, set
\[
U_{\mathrm{budget}}=24(d+3),\qquad
V_{\mathrm{budget}}=24\log(8(J+1)),\qquad
\delta_{\mathrm{count}}
=\frac{2\beta}{2\beta+D_{\gamma_{\min}}}>0,
\qquad
c_{\mathrm{budget}}=(2A_{\mathrm{mesh}})^{-2\beta}>0.
\]
The elementary bound
\[
\log n\le\frac{2}{\delta_{\mathrm{count}}}
n^{\delta_{\mathrm{count}}/2}
\]
and divergence of the positive powers allow one common starting sample
size for which
\[
U_{\mathrm{budget}}\log n
\le\frac{c_{\mathrm{budget}}}{2}n^{\delta_{\mathrm{count}}},
\qquad
V_{\mathrm{budget}}
\le\frac{c_{\mathrm{budget}}}{2}n^{\delta_{\mathrm{count}}}.
\]
Together with the endpoint and candidate bounds, these give
\[
\begin{aligned}
24s_n
&\le U_{\mathrm{budget}}\log n+V_{\mathrm{budget}}\\
&\le(2A_{\mathrm{mesh}}h_{\gamma_{\min},n})^{-2\beta}\\
&\le(2A_{\mathrm{mesh}}h_{\gamma,n})^{-2\beta}
\le h_{\mathrm{o}}^{-2\beta}.
\end{aligned}
\]
Fix \(n_{\mathrm{count}}\ge2\) large enough that all these bounds hold
uniformly for \(\gamma\in\Gamma\).

Finally, the oracle balance identity is
\[
nh_{\gamma,n}^{2\beta+D_\gamma}=1.
\]
Thus every microcell at \(h_{\mathrm{o}}\) satisfies
\[
\begin{aligned}
nq_{Q\ell}h_{\mathrm{o}}^{2\beta}
&\ge n\kappa_*h_{\mathrm{o}}^{2\beta+D_\gamma}\\
&\ge
\kappa_*A_{\mathrm{mesh}}^{2\beta+D_\gamma}
nh_{\gamma,n}^{2\beta+D_\gamma}
\ge4,
\end{aligned}
\]
or equivalently \(nq_{Q\ell}\ge4h_{\mathrm{o}}^{-2\beta}\).
% lean: CausalSmith.Stat.WeakOverlap.selectedCountEvent_uniform

\item Under \cref{obj:ass:iid}, a treated microcell count is a sum of
independent Bernoulli indicators. Each indicator has mean \(q_{Q\ell}\),
its centered absolute value is at most one, and its variance is
\(q_{Q\ell}(1-q_{Q\ell})\le q_{Q\ell}\). The Bernstein bound for these
bounded centered summands gives, for every \(u>0\),
\[
P^{\otimes n}\{|N_{Q\ell}-nq_{Q\ell}|\ge u\}
\le
2\exp\!\left\{-\frac{u^2}{2(2nq_{Q\ell}+u)}\right\}.
\]
For \(u=nq_{Q\ell}/2+4s_n\), the numerical inequality
\[
2s_n\left(2nq_{Q\ell}+u\right)\le u^2
\]
follows by expanding both sides. Hence
\[
P^{\otimes n}\left\{
|N_{Q\ell}-nq_{Q\ell}|
\ge\frac{nq_{Q\ell}}2+4s_n
\right\}\le2e^{-s_n}.
\]
Define
\[
\mathcal E_n^{\mathrm{count}}
=
\bigcap_{h\in\mathcal H_n}
\bigcap_{Q\in\mathcal Q_h}
\bigcap_{\ell=1}^{J}
\left\{
|N_{Q\ell}-nq_{Q\ell}|
<\frac{nq_{Q\ell}}2+4s_n
\right\}.
\]
This finite intersection is measurable. The union bound gives
\[
P^{\otimes n}\!\left((\mathcal E_n^{\mathrm{count}})^c\right)
\le2M_{\mathrm{count},n}e^{-s_n}
=
\frac{M_{\mathrm{count},n}}{M_{\mathrm{count},n}+1}\,n^{-2}
\le n^{-2}.
\]

On this event, the candidate from the preceding step satisfies
\[
N_{Q\ell}
\ge\frac{nq_{Q\ell}}2-4s_n
\ge2h_{\mathrm{o}}^{-2\beta}-4s_n
\ge h_{\mathrm{o}}^{-2\beta}.
\]
It is therefore feasible. The selector in \cref{obj:def:mesh-selector}
then gives
\[
0<\widehat h\le h_{\mathrm{o}}
\le2A_{\mathrm{mesh}}h_{\gamma,n}.
\]
At the selected mesh, feasibility and the budget imply
\[
N_{Q\ell}\ge\widehat h^{-2\beta}
\ge h_{\mathrm{o}}^{-2\beta}\ge24s_n.
\]
Using the upper deviation inequality first gives
\[
nq_{Q\ell}
\ge\frac23(N_{Q\ell}-4s_n)
\ge\frac59N_{Q\ell}
\ge\frac{40}{3}s_n.
\]
The lower deviation inequality now yields
\[
N_{Q\ell}
\ge\frac{nq_{Q\ell}}2-4s_n
\ge\frac{nq_{Q\ell}}5.
\]
These arguments apply to every partition cube, with the boundary
convention of \cref{obj:synth_3}.
% lean: CausalSmith.Stat.WeakOverlap.selectedCountEvent_of_oracle_witness

\item We next obtain the actual conditional coefficient bound. Define
the sampled design and its conditional sample law by
\[
\mathcal G_n=\sigma((X_i,A_i)_{i=1}^n),
\qquad
\nu_\omega
=
\mathcal L\!\left(\mathcal D_n\mid
(X_i,A_i)_{i=1}^n=(X_i(\omega),A_i(\omega))_{i=1}^n\right).
\]
Under the iid product law, disintegration coordinate by coordinate gives
\[
\mathcal L_{\nu_\omega}((Y_i)_{i=1}^n)
=
\bigotimes_{i=1}^n
\mathcal L_P(Y\mid X=X_i(\omega),A=A_i(\omega))
\]
for almost every \(\omega\). To see the product structure, conditional
probabilities of response rectangles factor into the product of the
coordinate conditional probabilities; these rectangles determine the
conditional response law.

For such a design, define the treated residuals on a conditional sample
\(\xi\) by
\[
\varepsilon_i^\omega(\xi)
=
\begin{cases}
Y_i(\xi)-\overline y_P(X_i(\omega)),&A_i(\omega)=1,\\
0,&A_i(\omega)=0.
\end{cases}
\]
They are independent under \(\nu_\omega\), and the treated conditional
bound from the first step, together with the zero residual in the second
case, gives
\[
\mathbb E_{\nu_\omega}\varepsilon_i^\omega=0,\qquad
\mathbb E_{\nu_\omega}e^{t\varepsilon_i^\omega}
\le e^{B^2t^2/2}\quad(t\in\mathbb R),
\]
with integrable exponential moments.

Fix a feasible candidate \(h\), a cube \(Q\in\mathcal Q_h\), and
\(\alpha\in\mathcal I_m\). Define
\[
x_{iQ}=\frac{X_i(\omega)-a_Q}{h},
\qquad
w_{iQ\alpha}
=
\frac1J\sum_{\ell=1}^{J}
\frac{\mathbf1\{A_i(\omega)=1,\ X_i(\omega)\in E_{Q\ell}\}}
     {N_{Q\ell}(\omega)}
[\widehat G_Q(\omega)^{-1}U(x_{iQ})]_\alpha.
\]
The equal-cell formula in \cref{obj:synth_4} is understood with
\(0^{-1}=0\), and with the zero matrix assigned as the inverse of a
singular Gram matrix. On the feasible meshes considered here,
\cref{obj:lem:template-conditioning} supplies the ordinary inverse and
\[
\|\widehat G_Q^{-1}\|_{\mathrm{op}}\le\frac2{\lambda_0},
\qquad \lambda_0>0.
\]

The conditional sample has the same design as \(\omega\) almost surely.
Expanding the coefficient formula therefore gives
\[
\widehat c_{Q,\alpha}(\xi)
=
\sum_{i=1}^n w_{iQ\alpha}
\mathbf1\{A_i(\omega)=1\}\overline y_P(X_i(\omega))
+
\sum_{i=1}^n w_{iQ\alpha}\varepsilon_i^\omega(\xi).
\]
The residuals are integrable and have mean zero, so integration proves
\[
\mathbb E_{\nu_\omega}\widehat c_{Q,\alpha}
=
\sum_{i=1}^n w_{iQ\alpha}
\mathbf1\{A_i(\omega)=1\}\overline y_P(X_i(\omega)).
\]
In particular, defining
\[
Z_{Q,\alpha}^\omega
=
\widehat c_{Q,\alpha}
-\mathbb E_{\nu_\omega}\widehat c_{Q,\alpha},
\qquad
Z_Q^\omega=(Z_{Q,\alpha}^\omega)_{\alpha\in\mathcal I_m},
\]
we have the conditional identity and moment bound
\[
Z_{Q,\alpha}^\omega
=
\sum_{i=1}^n w_{iQ\alpha}\varepsilon_i^\omega,
\qquad
\mathbb E_{\nu_\omega}e^{tZ_{Q,\alpha}^\omega}
=
\prod_{i=1}^n
\mathbb E_{\nu_\omega}e^{tw_{iQ\alpha}\varepsilon_i^\omega}
\le
\exp\!\left\{
\frac{B^2t^2}{2}\sum_{i=1}^n w_{iQ\alpha}^2
\right\}.
\]
Independence and coordinate exponential integrability also show that
every exponential moment in this display is integrable.

The normalized covariates appearing in a microcell lie in
\([0,1]^d\). Since every coordinate of \(U\) has absolute value at most
one there,
\[
\|U(x_{iQ})\|_2\le\sqrt{p_{\mathrm{coef}}}.
\]
Consequently, if
\[
R_{\mathrm{row}}
=
\frac{2\sqrt{p_{\mathrm{coef}}}}{J\lambda_0},
\]
then
\[
\left|
\frac1J[\widehat G_Q^{-1}U(x_{iQ})]_\alpha
\right|\le R_{\mathrm{row}}
\]
for observations counted in a microcell. Applying
\((\sum_{\ell=1}^{J}z_\ell)^2\le J\sum_{\ell=1}^{J}z_\ell^2\)
to the cell contributions to each weight, and then summing over the
observations, gives
\[
\begin{aligned}
\sum_{i=1}^n w_{iQ\alpha}^2
&\le
J\sum_{\ell=1}^{J}
\sum_{\substack{i:A_i=1\\X_i\in E_{Q\ell}}}
\frac{R_{\mathrm{row}}^2}{N_{Q\ell}^2}\\
&=
JR_{\mathrm{row}}^2\sum_{\ell=1}^{J}N_{Q\ell}^{-1}.
\end{aligned}
\]

For the mass comparison, define, at the fixed mesh,
\[
\mathcal S_Q=\{Q'\in\mathcal Q_h:q_{Q'}\le q_Q\},
\qquad
\rho_Q=|\mathcal S_Q|\ge1,
\]
and select weakest microcells by
\[
\ell_*(Q')\in
\operatorname*{arg\,min}_{1\le\ell\le J}q_{Q'\ell},
\qquad
\mathcal U_Q=\bigcup_{Q'\in\mathcal S_Q}E_{Q',\ell_*(Q')}.
\]
The selected microcells in different cubes are disjoint. Their union
has covariate mass at least \(\rho_Qc_f\eta^dh^d\), whereas its treated
mass is at most \(\rho_Qq_Q\). The setwise inequality of
\cref{obj:lem:ordered-mass} thus gives
\[
\rho_Qq_Q
\ge
\frac{\gamma-1}{\gamma}C^{-1/(\gamma-1)}
(\rho_Qc_f\eta^dh^d)^{\gamma/(\gamma-1)}.
\]
Dividing by \(\rho_Q>0\) and collecting powers yields the explicit
mass-rank bound
\[
q_Q\ge
c_{\mathrm{mass}}(\gamma)h^{D_\gamma}
\rho_Q^{1/(\gamma-1)}.
\]
At a feasible mesh satisfying the count lower bound from the preceding
step, each inverse count consequently satisfies
\[
N_{Q\ell}^{-1}
\le
\min\!\left\{
h^{2\beta},
\frac{5}{
n c_{\mathrm{mass}}(\gamma)h^{D_\gamma}
\rho_Q^{1/(\gamma-1)}}
\right\}.
\]
Define the fixed template factor
\[
A_{\mathrm{coef}}
=
B^2J^2R_{\mathrm{row}}^2
=
\frac{4B^2p_{\mathrm{coef}}}{\lambda_0^2}.
\]
The conditional coefficient bound is therefore
\[
\mathbb E_{\nu_\omega}e^{tZ_{Q,\alpha}^\omega}
\le
\exp\!\left[
\frac{A_{\mathrm{coef}}t^2}{2}
\min\!\left\{
h^{2\beta},
\frac{5}{
n c_{\mathrm{mass}}(\gamma)h^{D_\gamma}
\rho_Q^{1/(\gamma-1)}}
\right\}
\right].
\]

Choose a measurable full-measure set
\(\mathcal E_n^{\mathrm{cond}}\) on which the conditional identities,
exponential integrability, and bounds just established hold
simultaneously for all feasible candidate indices and coefficient
coordinates. Define
\[
\mathcal E_n
=
\mathcal E_n^{\mathrm{count}}\cap\mathcal E_n^{\mathrm{cond}},
\qquad
P^{\otimes n}(\mathcal E_n^{\mathrm{cond}})=1.
\]
Intersecting with this full-measure set preserves the failure bound:
\[
P^{\otimes n}(\mathcal E_n^c)\le n^{-2}.
\]
All subsequent conditional assertions hold for every \(\omega\in\mathcal
E_n\), at the selected mesh of that sample.
% lean: CausalSmith.Stat.WeakOverlap.selectedMesh_count_and_raw_mgf

\item We normalize the two variance scales uniformly. Define
\[
F_{\mathrm{mass}}=\max\{1,5/\kappa_*\}.
\]
For \(x\ge0\), \(y>0\), and \(c\ge\kappa_*\), we have
\[
\frac5{cy}\le F_{\mathrm{mass}}y^{-1},
\qquad
\min\{x,5/(cy)\}
\le F_{\mathrm{mass}}\min\{x,y^{-1}\}.
\]
For the second inequality, if \(x\le y^{-1}\), the left side is at most
\(x\le F_{\mathrm{mass}}x\); if \(x>y^{-1}\), it is at most
\(5/(cy)\le F_{\mathrm{mass}}y^{-1}\).
Applying this with the displayed mass coefficient and collecting
inverse powers gives
\[
\begin{aligned}
&A_{\mathrm{coef}}
\min\!\left\{
\widehat h^{2\beta},
\frac5{
n c_{\mathrm{mass}}(\gamma)\widehat h^{D_\gamma}
\rho_Q^{1/(\gamma-1)}}
\right\}\\
&\qquad\le
A_{\mathrm{coef}}F_{\mathrm{mass}}
\min\!\left\{
\widehat h^{2\beta},
n^{-1}\widehat h^{-D_\gamma}\rho_Q^{-1/(\gamma-1)}
\right\}.
\end{aligned}
\]
Set
\[
K_{\mathrm{count}}
=
\max\{2A_{\mathrm{mesh}},A_{\mathrm{coef}}F_{\mathrm{mass}}\}+1>0.
\]
The selected mesh and conditional moments thus obey
\[
0<\widehat h\le K_{\mathrm{count}}h_{\gamma,n},
\qquad
\mathbb E_{\nu_\omega}e^{tZ_{Q,\alpha}^\omega}
\le
\exp\!\left\{
\frac{K_{\mathrm{count}}t^2}{2}
\min\!\left(
\widehat h^{2\beta},
n^{-1}\widehat h^{-D_\gamma}\rho_Q^{-1/(\gamma-1)}
\right)
\right\}.
\]
If the cubes are enumerated in nondecreasing mass order, the first \(k\)
cubes all belong to \(\mathcal S_{Q_{(k)}}\). Hence
\[
\rho_{Q_{(k)}}\ge k,
\qquad
\rho_{Q_{(k)}}^{-1/(\gamma-1)}
\le k^{-1/(\gamma-1)}.
\]
This converts the displayed bound to the rank-\(k\) envelope, including
ties.
% lean: CausalSmith.Stat.WeakOverlap.selectedMesh_count_and_rank_mgf

\item We establish the scalar maximal estimate used to control the
coefficient vectors. Begin with a clipped Gaussian integral. For an
integer \(R_{\mathrm{cut}}\ge1\) and a scale \(a>0\), define
\[
s_{\mathrm{clip}}=\sqrt{1+\log R_{\mathrm{cut}}},
\qquad
r_{\mathrm{clip}}=a s_{\mathrm{clip}},
\qquad
g_{\mathrm{pre}}(u)
=
\min\!\left\{1,\,
2R_{\mathrm{cut}}e^{-u^2/(2a^2)}\right\}.
\]
We have \(s_{\mathrm{clip}}^2\ge1\) and
\(\log(2R_{\mathrm{cut}})\le s_{\mathrm{clip}}^2\).
For \(u\le4r_{\mathrm{clip}}\),
\[
g_{\mathrm{pre}}(u)\le1\le e^{4-u/r_{\mathrm{clip}}}.
\]
For \(u>4r_{\mathrm{clip}}\), put
\[
x_{\mathrm{clip}}=u/r_{\mathrm{clip}}>4.
\]
The nonnegativity of
\[
(s_{\mathrm{clip}}^2-1)(x_{\mathrm{clip}}^2/2-1)
\]
and the elementary quadratic inequality
\(1-x_{\mathrm{clip}}^2/2\le4-x_{\mathrm{clip}}\) give
\[
s_{\mathrm{clip}}^2
-s_{\mathrm{clip}}^2x_{\mathrm{clip}}^2/2
\le4-x_{\mathrm{clip}}.
\]
Therefore
\[
g_{\mathrm{pre}}(u)
\le
\exp\!\left\{
\log(2R_{\mathrm{cut}})
-s_{\mathrm{clip}}^2x_{\mathrm{clip}}^2/2
\right\}
\le e^{4-u/r_{\mathrm{clip}}}.
\]
Integrating this exponential majorant proves
\[
\int_0^\infty g_{\mathrm{pre}}(u)\,du
\le e^4a\sqrt{1+\log R_{\mathrm{cut}}}.
\]
% lean: Causalean.Mathlib.Probability.SubGaussian.clipped_gaussian_integral

\item Fix any exponent \(\alpha_{\mathrm{ord}}>0\). For an integer
\(N\ge1\), retain \(R_{\mathrm{cut}}\ge1\) and define the clipped suffix
envelope
\[
g_{\mathrm{suf}}(u)
=
\min\!\left\{1,\,
\sum_{\substack{1\le k\le N\\k\ge R_{\mathrm{cut}}}}
2\exp\!\left[
-\frac{u^2}{
2a^2(R_{\mathrm{cut}}/k)^{\alpha_{\mathrm{ord}}}}
\right]\right\}.
\]
Partition its indices into dyadic blocks by
\[
L_j=2^jR_{\mathrm{cut}},\qquad
\mathcal B_j
=
\{k:1\le k\le N,\ L_j\le k<2L_j\},
\qquad j\in\mathbb N_0,
\]
and define
\[
g_j(u)
=
\min\!\left\{1,\,
\sum_{k\in\mathcal B_j}
2\exp\!\left[
-\frac{u^2}{
2a^2(R_{\mathrm{cut}}/k)^{\alpha_{\mathrm{ord}}}}
\right]\right\}.
\]
Each block has at most \(L_j\) indices, and within it
\[
a^2(R_{\mathrm{cut}}/k)^{\alpha_{\mathrm{ord}}}
\le a^2(R_{\mathrm{cut}}/L_j)^{\alpha_{\mathrm{ord}}}.
\]
The preceding clipped Gaussian calculation, at scale
\(a(R_{\mathrm{cut}}/L_j)^{\alpha_{\mathrm{ord}}/2}\), consequently gives
\[
\int_0^\infty g_j(u)\,du
\le
e^4a\,2^{-j\alpha_{\mathrm{ord}}/2}
\sqrt{1+\log(2^jR_{\mathrm{cut}})}.
\]
Since \(\log2\le1\) and \(\log R_{\mathrm{cut}}\ge0\),
\[
\sqrt{1+\log(2^jR_{\mathrm{cut}})}
\le(j+1)\sqrt{1+\log R_{\mathrm{cut}}}.
\]
Define
\[
q_{\mathrm{geom}}
=
e^{-\alpha_{\mathrm{ord}}\log2/2}\in(0,1),
\qquad
D_{\mathrm{geom}}(\alpha_{\mathrm{ord}})
=
1+\sum_{j=0}^{\infty}(j+1)q_{\mathrm{geom}}^j<\infty.
\]
Convergence follows from the geometric series and its linearly weighted
series.

The dyadic blocks partition the suffix. For nonnegative summands,
clipping a sum at one is bounded by the sum of the individually clipped
block sums. Only finitely many blocks meet \(\{1,\ldots,N\}\), so
integration gives
\[
\begin{aligned}
\int_0^\infty g_{\mathrm{suf}}(u)\,du
&\le\sum_{j=0}^{N}\int_0^\infty g_j(u)\,du\\
&\le
e^4a\sqrt{1+\log R_{\mathrm{cut}}}
\sum_{j=0}^{N}(j+1)q_{\mathrm{geom}}^j\\
&\le
e^4D_{\mathrm{geom}}(\alpha_{\mathrm{ord}})
a\sqrt{1+\log R_{\mathrm{cut}}}.
\end{aligned}
\]
The finite Gaussian sums and their clipped versions are integrable, so
all these integrations are legitimate.
% lean: Causalean.Mathlib.Probability.SubGaussian.decaying_gaussian_suffix_integral

\item Now let \(\nu\) be a probability law, let \(a,b>0\), and let
\(Z_1,\ldots,Z_N\) have the exponential integrability and moment envelope
with exponent \(\alpha_{\mathrm{ord}}\). Define
\[
W_N=\max\bigl(\{0\}\cup\{|Z_k|:1\le k\le N\}\bigr).
\]
For \(N=0\), this definition gives \(W_N=0\), and the maximal inequality
follows immediately. Suppose \(N\ge1\), and define
\[
v_k=\min\{a^2,b^2k^{-\alpha_{\mathrm{ord}}}\}>0,
\qquad
\tau_{\mathrm{tail}}(u)=\min\!\left\{1,\sum_{k=1}^{N}2e^{-u^2/(2v_k)}\right\}.
\]
Exponential Markov bounds for \(Z_k\) and \(-Z_k\), optimized at
\(u/v_k\), give
\[
\nu\{|Z_k|>u\}\le2e^{-u^2/(2v_k)}
\quad(u\ge0).
\]
The finite union bound, together with the probability bound by one,
therefore yields
\[
\nu\{W_N>u\}\le\tau_{\mathrm{tail}}(u).
\]
Exponential moments at \(1\) and \(-1\) imply integrability of every
\(|Z_k|\), and hence of \(W_N\). The layer-cake identity gives
\[
\mathbb E_\nu W_N
=
\int_0^\infty\nu\{W_N>u\}\,du
\le\int_0^\infty\tau_{\mathrm{tail}}(u)\,du.
\]

To bound the last integral, set
\[
R_{\mathrm{ratio}}=(b/a)^{2/\alpha_{\mathrm{ord}}}>0,
\qquad
R_{\mathrm{cut}}=\lceil R_{\mathrm{ratio}}\rceil\ge1.
\]
Then
\[
R_{\mathrm{ratio}}\le R_{\mathrm{cut}}
\le2\max\{1,R_{\mathrm{ratio}}\},
\qquad
1+\log R_{\mathrm{cut}}
\le2(1+\log_+(R_{\mathrm{ratio}})).
\]
The latter bound follows by taking logarithms of the preceding
inequality, using
\(\log\max\{1,R_{\mathrm{ratio}}\}
=\log_+(R_{\mathrm{ratio}})\) and \(\log2\le1\).
It implies
\[
\sqrt{1+\log R_{\mathrm{cut}}}
\le2\sqrt{1+\log_+(R_{\mathrm{ratio}})}.
\]
Furthermore,
\[
R_{\mathrm{ratio}}^{\alpha_{\mathrm{ord}}}=(b/a)^2,
\qquad
b^2\le a^2R_{\mathrm{cut}}^{\alpha_{\mathrm{ord}}},
\qquad
v_k\le a^2(R_{\mathrm{cut}}/k)^{\alpha_{\mathrm{ord}}}.
\]
Split the sum defining \(\tau_{\mathrm{tail}}\) at \(k<R_{\mathrm{cut}}\) and
\(k\ge R_{\mathrm{cut}}\). The prefix has at most \(R_{\mathrm{cut}}\)
terms, each with variance at most \(a^2\). The suffix has the preceding
decaying variance bound. Since
\[
\min\{1,x+y\}\le\min\{1,x\}+\min\{1,y\}
\qquad(x,y\ge0),
\]
we obtain
\[
\tau_{\mathrm{tail}}(u)\le g_{\mathrm{pre}}(u)+g_{\mathrm{suf}}(u).
\]
The two integral estimates prove
\[
\mathbb E_\nu W_N
\le
K_{\mathrm{sc}}(\alpha_{\mathrm{ord}})
a\sqrt{1+\log_+((b/a)^{2/\alpha_{\mathrm{ord}}})},
\]
where the positive constant may be taken as
\[
K_{\mathrm{sc}}(\alpha_{\mathrm{ord}})
=
2e^4\bigl(1+D_{\mathrm{geom}}(\alpha_{\mathrm{ord}})\bigr).
\]
% lean: Causalean.Mathlib.Probability.SubGaussian.orderedSubGaussianMaximal

\item To obtain a constant uniform over the adaptation interval, define
\[
\alpha_{\min}=\frac1{\gamma_{\max}-1}>0,\qquad
\alpha_{\max}=\frac1{\gamma_{\min}-1},
\qquad
R_{\mathrm{exp}}=\frac{\alpha_{\max}}{\alpha_{\min}}\ge1.
\]
Taking reciprocals of positive denominators shows
\[
\alpha_{\min}\le\frac1{\gamma-1}\le\alpha_{\max}
\qquad(\gamma\in\Gamma).
\]

For any \(\alpha_{\mathrm{ord}}\in[\alpha_{\min},\alpha_{\max}]\) and
\(k\ge1\),
\[
k^{-\alpha_{\mathrm{ord}}}\le k^{-\alpha_{\min}}.
\]
Thus the scalar estimate at \(\alpha_{\min}\) applies to the envelope
with exponent \(\alpha_{\mathrm{ord}}\). Its logarithmic factor satisfies
\[
1+\log_+((b/a)^{2/\alpha_{\min}})
\le
R_{\mathrm{exp}}
\bigl(1+\log_+((b/a)^{2/\alpha_{\mathrm{ord}}})\bigr).
\]
Indeed, if \(\log(b/a)\ge0\), use
\[
\frac2{\alpha_{\min}}
\le R_{\mathrm{exp}}\frac2{\alpha_{\mathrm{ord}}};
\]
if \(\log(b/a)<0\), both positive-part logarithms vanish and the
inequality is \(1\le R_{\mathrm{exp}}\).
Taking square roots establishes a uniform scalar constant
\[
M_\Gamma
=
K_{\mathrm{sc}}(\alpha_{\min})\sqrt{R_{\mathrm{exp}}}>0.
\]

For a finite nonempty family of vectors with coordinates in
\(\mathcal I_m\), the pointwise inequality is
\[
\max_Q\|Z_Q\|_\infty
\le
\sum_{\alpha\in\mathcal I_m}\max_Q|Z_{Q,\alpha}|.
\]
Coordinate exponential integrability implies integrability of these
finite maxima and of the vector maximum. For example,
\[
|Z_{Q,\alpha}|
\le e^{Z_{Q,\alpha}}+e^{-Z_{Q,\alpha}},
\]
and each finite maximum is bounded by the sum of the corresponding
absolute values. Applying the uniform scalar estimate coordinatewise
and integrating the pointwise inequality gives
\[
\mathbb E_\nu\max_Q\|Z_Q\|_\infty
\le
p_{\mathrm{coef}}M_\Gamma a
\sqrt{1+\log_+((b/a)^{2/\alpha_{\mathrm{ord}}})}.
\]
% lean: CausalSmith.Stat.WeakOverlap.orderedSubGaussianMaximal_vector_uniform_of_exp

\item The preceding vector estimate applies to mass-ranked cube
envelopes. Enumerate the cubes in nondecreasing \(q_Q\), breaking ties
by any fixed ordering of the finite cube family. This enumeration is a
bijection, so
\[
\max_Q\|Z_Q\|_\infty
=
\max_{1\le k\le|\mathcal Q_h|}\|Z_{Q_{(k)}}\|_\infty.
\]
As established above, \(\rho_{Q_{(k)}}\ge k\). For any nonnegative
variance factor \(y\),
\[
\min\{x,y\rho_{Q_{(k)}}^{-\alpha_{\mathrm{ord}}}\}
\le
\min\{x,yk^{-\alpha_{\mathrm{ord}}}\}.
\]
Therefore an envelope with variance proxy
\(K_{\mathrm{count}}\min\{x,y\rho_Q^{-\alpha_{\mathrm{ord}}}\}\)
has the ordered envelope required by the vector estimate whenever
\[
a^2=K_{\mathrm{count}}x,\qquad
b^2=K_{\mathrm{count}}y.
\]
The cube family is nonempty at every dyadic mesh, so the vector estimate
applies throughout the selected-mesh argument.
% lean: CausalSmith.Stat.WeakOverlap.selectedMesh_cube_rank_vector_maximal

\item We now fix the final constants. Define
\[
E_\Gamma
=
\max\{1,\,2\beta(\gamma_{\max}-1)+d\gamma_{\max}\},
\]
\[
K=
\max\!\left\{
K_{\mathrm{count}},
p_{\mathrm{coef}}M_\Gamma
\sqrt{K_{\mathrm{count}}}\sqrt{E_\Gamma}
\right\}+1,
\qquad
K'=K\max\{1,K^\beta\},
\qquad
n_0=n_{\mathrm{count}}.
\]
These constants are positive and depend only on the fixed parameters
and interval endpoints.

Fix \(n\ge n_0\), \(\gamma\in\Gamma\), \(P\in\mathcal P_\gamma\), and
\(\omega\in\mathcal E_n\). Conditional on its sampled design, keep
\(\widehat h=\widehat h(\omega)\) fixed and define the scales
\[
a_{\mathrm{sel}}
=
\sqrt{K_{\mathrm{count}}}\widehat h^\beta,
\qquad
b_{\mathrm{sel}}
=
\sqrt{K_{\mathrm{count}}}\,
n^{-1/2}\widehat h^{-D_\gamma/2},
\qquad
\alpha_{\mathrm{sel}}=\frac1{\gamma-1}.
\]
All three are positive, and
\[
a_{\mathrm{sel}}^2
=K_{\mathrm{count}}\widehat h^{2\beta},
\qquad
b_{\mathrm{sel}}^2
=K_{\mathrm{count}}n^{-1}\widehat h^{-D_\gamma},
\qquad
\alpha_{\mathrm{sel}}\in[\alpha_{\min},\alpha_{\max}].
\]
Applying the ranked vector estimate under \(\nu_\omega\) gives
\[
\mathbb E_{\nu_\omega}
\max_{Q\in\mathcal Q_{\widehat h}}\|Z_Q^\omega\|_\infty
\le
p_{\mathrm{coef}}M_\Gamma a_{\mathrm{sel}}
\sqrt{
1+\log_+\bigl((b_{\mathrm{sel}}/a_{\mathrm{sel}})
^{2/\alpha_{\mathrm{sel}}}\bigr)
}.
\]
The same exponential integrability supplies integrability of the
maximum on the left.
% lean: CausalSmith.Stat.WeakOverlap.selectedMesh_envelope

\item The oracle definition and \(2\beta+D_\gamma>0\) imply
\[
n^{-1/2}
=h_{\gamma,n}^{\,\beta+D_\gamma/2}.
\]
After cancelling \(\sqrt{K_{\mathrm{count}}}\), this gives
\[
\frac{b_{\mathrm{sel}}}{a_{\mathrm{sel}}}
=
\left(\frac{h_{\gamma,n}}{\widehat h}\right)^{
\beta+D_\gamma/2}.
\]
Define the resulting logarithmic exponent by
\[
e_{\mathrm{sel}}
=
\left(\beta+\frac{D_\gamma}{2}\right)
\frac2{\alpha_{\mathrm{sel}}}
=
2\beta(\gamma-1)+d\gamma.
\]
Its reciprocal dependence cancels, and
\[
0\le e_{\mathrm{sel}}
\le2\beta(\gamma_{\max}-1)+d\gamma_{\max}
\le E_\Gamma.
\]

For a positive ratio \(u\) and an exponent \(\zeta_{\mathrm{exp}}\ge0\),
\[
1+\log_+(u^{\zeta_{\mathrm{exp}}})
\le\max\{1,\zeta_{\mathrm{exp}}\}(1+\log_+(u)).
\]
If \(\log u\ge0\), this follows from
\(1+\zeta_{\mathrm{exp}}\log u\le\max\{1,\zeta_{\mathrm{exp}}\}(1+\log u)\).
If \(\log u<0\), both positive parts vanish and the assertion is
\(1\le\max\{1,\zeta_{\mathrm{exp}}\}\). Applying these two cases to the oracle ratio proves
\[
\sqrt{
1+\log_+\bigl((b_{\mathrm{sel}}/a_{\mathrm{sel}})
^{2/\alpha_{\mathrm{sel}}}\bigr)
}
\le
\sqrt{E_\Gamma}\,
\sqrt{1+\log_+(h_{\gamma,n}/\widehat h)}.
\]
% lean: CausalSmith.Stat.WeakOverlap.selectedMesh_maximal_oracle_log_ratio

\item Combining the preceding estimates and the definition of \(K\)
yields
\[
\begin{aligned}
\mathbb E_{\nu_\omega}
\max_{Q\in\mathcal Q_{\widehat h}}\|Z_Q^\omega\|_\infty
&\le
p_{\mathrm{coef}}M_\Gamma
\sqrt{K_{\mathrm{count}}}\sqrt{E_\Gamma}\,
\widehat h^\beta
\sqrt{1+\log_+(h_{\gamma,n}/\widehat h)}\\
&\le
K\widehat h^\beta
\sqrt{1+\log_+(h_{\gamma,n}/\widehat h)}.
\end{aligned}
\]
Since \(K\ge K_{\mathrm{count}}\), the mesh bound also becomes
\[
0<\widehat h\le Kh_{\gamma,n}.
\]
The variance proxies are nonnegative, so increasing
\(K_{\mathrm{count}}\) to \(K\) in the conditional exponential bound
preserves that inequality. Together with the rank comparison, this
gives precisely
\[
\mathbb E_{\nu_\omega}e^{tZ_{Q_{(k)},\alpha}^\omega}
\le
\exp\!\left\{
\frac{Kt^2}{2}
\min\!\left(
\widehat h^{2\beta},
n^{-1}\widehat h^{-D_\gamma}k^{-1/(\gamma-1)}
\right)
\right\}.
\]
The count bound and all exponential integrability assertions are
retained on the same event \(\mathcal E_n\).
% lean: CausalSmith.Stat.WeakOverlap.selectedMesh_envelope

\item For the deterministic comparison with the oracle rate, define
\[
s_{\mathrm{mesh}}=\frac{\widehat h}{h_{\gamma,n}},
\qquad
0<s_{\mathrm{mesh}}\le K.
\]
If \(s_{\mathrm{mesh}}\le1\), then
\[
0\le\log(1/s_{\mathrm{mesh}})
\le1/s_{\mathrm{mesh}}-1,
\]
and consequently
\[
\sqrt{1+\log_+(1/s_{\mathrm{mesh}})}
\le\sqrt{1/s_{\mathrm{mesh}}}
\le1/s_{\mathrm{mesh}}.
\]
Since \(\beta>1\),
\[
s_{\mathrm{mesh}}^\beta\le s_{\mathrm{mesh}},
\qquad
s_{\mathrm{mesh}}^\beta
\sqrt{1+\log_+(1/s_{\mathrm{mesh}})}
\le1.
\]
If \(s_{\mathrm{mesh}}>1\), the positive-part logarithm vanishes, and
\[
s_{\mathrm{mesh}}^\beta
\sqrt{1+\log_+(1/s_{\mathrm{mesh}})}
=s_{\mathrm{mesh}}^\beta\le K^\beta.
\]
The two cases give
\[
\begin{aligned}
\widehat h^\beta
\sqrt{1+\log_+(h_{\gamma,n}/\widehat h)}
&=
h_{\gamma,n}^\beta s_{\mathrm{mesh}}^\beta
\sqrt{1+\log_+(1/s_{\mathrm{mesh}})}\\
&\le\max\{1,K^\beta\}r_{\gamma,n}.
\end{aligned}
\]
Multiplying by \(K\) proves the asserted bound with
\(K'=K\max\{1,K^\beta\}\).
% lean: CausalSmith.Stat.WeakOverlap.selectedMesh_oracle_rate_log_bound

\item Finally, for the additional maximal assertion, take
\(\nu=\mu\) and \(\alpha_{\mathrm{ord}}=\alpha>0\) in the scalar
calculation above, and set
\[
K_\alpha=K_{\mathrm{sc}}(\alpha)>0.
\]
Its hypotheses are exactly the stated exponential integrability and
moment envelope. Thus, including the case \(N=0\),
\[
\mathbb E_\mu
\max\bigl(\{0\}\cup\{|Z_k|:1\le k\le N\}\bigr)
\le
K_\alpha a
\sqrt{1+\max\{0,\log((b/a)^{2/\alpha})\}}.
\]
This completes all assertions.
% lean: Causalean.Mathlib.Probability.SubGaussian.orderedSubGaussianMaximal
\end{enumerate}
\end{proof}
